% GENERATED FILE. Edit canonical lessons, then run npm run generate:books.
% canonical-source-hash: fnv1a64:2f49057a16038032
% canonical-lessons: ZH-C164-daban, ZH-C164-xitou, ZH-C164-chuifengji, ZH-C164-huazhuang, ZH-C164-gualian

\chapter{Getting Ready}
\label{ch:zh-getting-ready}
\hlchaptermodality{164}

\begin{chapteropening}
\textbf{By the end of this chapter:} \emph{I can talk about dressing up, washing and drying my hair, make-up and shaving.}

Complete the last lesson of chapter 164: I can talk about dressing up, washing and drying my hair, make-up and shaving.
\end{chapteropening}

\section[dǎban]{dǎban — to dress up; to get dressed up}

\label{lesson:ZH-C164-daban}



\begin{quote}
Before the new one: say the Mandarin for next time, then the Mandarin for all together.
\end{quote}

\subsection*{You'll want to know: dǎban}
\emph{dǎban} — \textquotedblleft{}to dress up; to get dressed up\textquotedblright{}.

\begin{quote}
\emph{tā zài dǎban wǒmen děng yíhuìr} — She's getting ready; let's wait a moment.
\end{quote}

Its characters come later; for now it is a word for the ear.

\begin{cousinweb}
\emph{dǎban} is built from \zh{打} 'to strike' and \emph{ban}, both inherited from Middle Chinese, \emph{taengX} and \emph{peanH} (Guangyun, in Baxter's transcription), which glosses the second with this very word, 'to dress up'. Dǎ is not the sound \emph{taengX} would regularly give, and in the word the second loses its tone.
\end{cousinweb}

\begin{grammarlens}[title={What you've built}]
One more word for saying what I do.
\end{grammarlens}

\subsection*{Your turn}
\begin{itemize}
\raggedright
  \item \emph{Say it:} \emph{dǎban}
  \item \emph{Say it:} \emph{dǎban}, once more
  \item \emph{Say it:} \emph{tā zài dǎban wǒmen děng yíhuìr} — She's getting ready; let's wait a moment.
  \item \emph{Recall:} say the Mandarin for next time, then the Mandarin for all together, then say \emph{dǎban} again
\end{itemize}

\begin{tcolorbox}[breakable,colback=teal!4,colframe=teal!35!black,title={Before you move on}]
What does dǎban mean? (\textquotedblleft{}To dress up; to get dressed up\textquotedblright{}.) Say it once more.
\end{tcolorbox}
\section[xǐtóu]{xǐtóu — to wash one's hair}

\label{lesson:ZH-C164-xitou}



\begin{quote}
Before the new one: say the Mandarin for all together, then the Mandarin for to dress up.
\end{quote}

\subsection*{You'll want to know: xǐtóu}
\emph{xǐtóu} — \textquotedblleft{}to wash one's hair\textquotedblright{}.

\begin{quote}
\emph{wǒ měitiān zǎoshang xǐtóu} — I wash my hair every morning.
\end{quote}

Its characters come later; for now it is a word for the ear.

\begin{cousinweb}
\emph{xǐtóu} is built from \zh{洗} 'to wash' and \emph{tóu} 'head', both inherited from Middle Chinese, \emph{sejX} and \emph{duw} (Guangyun, in Baxter's transcription).
\end{cousinweb}

\begin{grammarlens}[title={What you've built}]
One more word for saying what I do.
\end{grammarlens}

\subsection*{Your turn}
\begin{itemize}
\raggedright
  \item \emph{Say it:} \emph{xǐtóu}
  \item \emph{Say it:} \emph{xǐtóu}, once more
  \item \emph{Say it:} \emph{wǒ měitiān zǎoshang xǐtóu} — I wash my hair every morning.
  \item \emph{Recall:} say the Mandarin for all together, then the Mandarin for to dress up, then say \emph{xǐtóu} again
\end{itemize}

\begin{tcolorbox}[breakable,colback=teal!4,colframe=teal!35!black,title={Before you move on}]
What does xǐtóu mean? (\textquotedblleft{}To wash one's hair\textquotedblright{}.) Say it once more.
\end{tcolorbox}
\section[chuīfēngjī]{chuīfēngjī — hair dryer}

\label{lesson:ZH-C164-chuifengji}



\begin{quote}
Before the new one: say the Mandarin for to dress up, then the Mandarin for to wash one's hair.
\end{quote}

\subsection*{You'll want to know: chuīfēngjī}
\emph{chuīfēngjī} — \textquotedblleft{}hair dryer\textquotedblright{}.

\begin{quote}
\emph{xǐtóu yǐhòu wǒ yòng chuīfēngjī} — After washing my hair I use the hair dryer.
\end{quote}

Its characters come later; for now it is a word for the ear.

\begin{cousinweb}
\emph{chuīfēngjī} is built from \emph{chuī} 'to blow', \emph{fēng} 'wind' and \zh{机} 'machine'. The last two are inherited from Middle Chinese, \emph{pjuwng} and \emph{kj+j} (Guangyun, in Baxter's transcription); the first is listed there as \emph{tsyhwe}, glossed 'to breathe out': the machine that blows wind.
\end{cousinweb}

\begin{grammarlens}[title={What you've built}]
One more word for saying what I do.
\end{grammarlens}

\subsection*{Your turn}
\begin{itemize}
\raggedright
  \item \emph{Say it:} \emph{chuīfēngjī}
  \item \emph{Say it:} \emph{chuīfēngjī}, once more
  \item \emph{Say it:} \emph{xǐtóu yǐhòu wǒ yòng chuīfēngjī} — After washing my hair I use the hair dryer.
  \item \emph{Recall:} say the Mandarin for to dress up, then the Mandarin for to wash one's hair, then say \emph{chuīfēngjī} again
\end{itemize}

\begin{tcolorbox}[breakable,colback=teal!4,colframe=teal!35!black,title={Before you move on}]
What does chuīfēngjī mean? (\textquotedblleft{}Hair dryer\textquotedblright{}.) Say it once more.
\end{tcolorbox}
\section[huàzhuāng]{huàzhuāng — to put on make-up}

\label{lesson:ZH-C164-huazhuang}



\begin{quote}
Before the new one: say the Mandarin for to wash one's hair, then the Mandarin for a hair dryer.
\end{quote}

\subsection*{You'll want to know: huàzhuāng}
\emph{huàzhuāng} — \textquotedblleft{}to put on make-up\textquotedblright{}.

\begin{quote}
\emph{wǒ māma měi cì chūmén yǐqián dōu huàzhuāng} — My mum puts on make-up every time before she goes out.
\end{quote}

Its characters come later; for now it is a word for the ear.

\begin{cousinweb}
\emph{huàzhuāng} is built from \emph{huà} 'to change' and \emph{zhuāng} 'to adorn', both inherited from Middle Chinese, \emph{xwaeH} and \emph{tsrjang} (Guangyun, in Baxter's transcription).
\end{cousinweb}

\begin{grammarlens}[title={What you've built}]
One more word for saying what I do.
\end{grammarlens}

\subsection*{Your turn}
\begin{itemize}
\raggedright
  \item \emph{Say it:} \emph{huàzhuāng}
  \item \emph{Say it:} \emph{huàzhuāng}, once more
  \item \emph{Say it:} \emph{wǒ māma měi cì chūmén yǐqián dōu huàzhuāng} — My mum puts on make-up every time before she goes out.
  \item \emph{Recall:} say the Mandarin for to wash one's hair, then the Mandarin for a hair dryer, then say \emph{huàzhuāng} again
\end{itemize}

\begin{tcolorbox}[breakable,colback=teal!4,colframe=teal!35!black,title={Before you move on}]
What does huàzhuāng mean? (\textquotedblleft{}To put on make-up\textquotedblright{}.) Say it once more.
\end{tcolorbox}
\section[guāliǎn]{guāliǎn — to have a shave; to shave one's face}

\label{lesson:ZH-C164-gualian}



\begin{quote}
Before the new one: say the Mandarin for a hair dryer, then the Mandarin for to put on make-up.
\end{quote}

\subsection*{You'll want to know: guāliǎn}
\emph{guāliǎn} — \textquotedblleft{}to have a shave; to shave one's face\textquotedblright{}.

\begin{quote}
\emph{bàba měitiān zǎoshang guāliǎn} — Dad shaves every morning.
\end{quote}

Its characters come later; for now it is a word for the ear.

\begin{cousinweb}
\emph{guāliǎn} is built from \emph{guā} 'to scrape' and \emph{liǎn} 'face'. The first is inherited from Middle Chinese \emph{kwaet} (Guangyun, in Baxter's transcription), glossed 'to scrape and pare'; the dictionary glosses the second, \emph{leamX}, as the name of a kind of stew, so how it came to mean face is not settled.
\end{cousinweb}

\begin{grammarlens}[title={What you've built}]
One more word for saying what I do.
\end{grammarlens}

\subsection*{Your turn}
\begin{itemize}
\raggedright
  \item \emph{Say it:} \emph{guāliǎn}
  \item \emph{Say it:} \emph{guāliǎn}, once more
  \item \emph{Say it:} \emph{bàba měitiān zǎoshang guāliǎn} — Dad shaves every morning.
  \item \emph{Recall:} say the Mandarin for a hair dryer, then the Mandarin for to put on make-up, then say \emph{guāliǎn} again
\end{itemize}

\begin{tcolorbox}[breakable,colback=teal!4,colframe=teal!35!black,title={Before you move on}]
What does guāliǎn mean? (\textquotedblleft{}To have a shave; to shave one's face\textquotedblright{}.) Say it once more.
\end{tcolorbox}
